\documentclass[11pt,a4paper]{article}

% --------------------
% Packages
% --------------------
\usepackage[a4paper,margin=2.5cm]{geometry}
\usepackage{amsmath}
\usepackage{graphicx}
\usepackage[
  colorlinks=true,
  linkcolor=black,
  urlcolor=blue,
  citecolor=black
]{hyperref}

\usepackage{setspace}
\usepackage{listings}
\usepackage{xcolor}

\setstretch{1.15}

% --------------------
% Code formatting
% --------------------
\lstset{
  basicstyle=\ttfamily\small,
  breaklines=true,
  frame=single,
  numbers=left,
  numberstyle=\tiny,
  keywordstyle=\color{blue},
  commentstyle=\color{gray}
}

% --------------------
% Title
% --------------------
\title{Planning Notes and Methods Development\\
Low-Cost Laser-Based Pipe Inspection System}

\author{Luis Nunes}
\date{\today}

\begin{document}

\maketitle

\tableofcontents
\newpage

% ==================================================
\section{Introduction}

This document records the evolving design decisions, software structure, and processing methods for the laser-based pipe inspection system.

It is intended as a technical planning and reference document to support later report writing and system refinement.

Entries are dated to track development progress.

% ==================================================
\section{System Overview}

The system consists of:

\begin{itemize}
  \item Raspberry Pi 5
  \item Raspberry Pi Camera (v2 / Camera Module 3 planned)
  \item Red line laser module (635--660 nm)
  \item Adjustable mechanical mount
  \item Python-based image processing pipeline
\end{itemize}

The system uses laser triangulation to estimate surface depth inside pipes.

% ==================================================
\section{Project Directory Structure}

\textbf{Date: 04 February 2026}

Current project structure:

\begin{lstlisting}
laser_project/
├── cfg/
│   └── config.yaml
├── src/
│   ├── utils.py
│   ├── extract_laser_line.py
│   ├── calibrate_camera.py
│   ├── calibrate_laser_plane.py
│   └── reconstruct_depth.py
├── data/
│   ├── raw/
│   └── processed/
└── calib/
\end{lstlisting}

This modular structure separates configuration, processing, calibration, and data.

% ==================================================
\section{Configuration System}

\textbf{Date: 04 February 2026}

System parameters are stored in a YAML configuration file.

Main sections:

\begin{itemize}
  \item \texttt{camera}: Intrinsic parameters
  \item \texttt{laser}: HSV detection thresholds
  \item \texttt{processing}: ROI and filtering parameters
\end{itemize}

Example structure:

\begin{lstlisting}
processing:
  roi: null
  blur_ksize: 3
  morph_ksize: 3
  min_pixels_per_row: 5
  subpixel: true
\end{lstlisting}

This allows rapid tuning without modifying source code.

% ==================================================
\section{Utility Functions}

\textbf{Date: 05 February 2026}

Common processing functions are stored in \texttt{utils.py}.

Main functions:

\begin{itemize}
  \item \texttt{load\_config}: Loads YAML configuration
  \item \texttt{crop\_roi}: Extracts region of interest
  \item \texttt{hsv\_mask\_red}: Detects red laser pixels
  \item \texttt{apply\_morph}: Cleans binary mask
\end{itemize}

This promotes code reuse and maintainability.

% ==================================================
\section{Laser Line Extraction Method}

\textbf{Date: 05 February 2026}

The current laser extraction method operates row-by-row.

Processing stages:

\begin{enumerate}
  \item Load input image
  \item Crop ROI (if enabled)
  \item Apply Gaussian blur
  \item Convert to HSV colour space
  \item Threshold red pixels
  \item Apply morphological filtering
  \item Extract centre per image row
\end{enumerate}

For each row $y$:

\begin{equation}
x(y) = \frac{1}{N}\sum_{i=1}^{N} x_i
\end{equation}

where $x_i$ are detected laser pixels.

This yields one centre point per row.

\subsection{Assumptions}

The method assumes:

\begin{itemize}
  \item A single dominant laser stripe
  \item Near-vertical stripe orientation
  \item One intersection per image row
\end{itemize}

Violations of these assumptions (e.g. reflections, rings) reduce accuracy.

% ==================================================
\section{Output Data}

\textbf{Date: 06 February 2026}

Each extraction run produces:

\begin{itemize}
  \item Binary mask image
  \item Overlay image with centreline
  \item CSV file of $(y,x)$ coordinates
\end{itemize}

Example CSV format:

\begin{lstlisting}
y,x
124,532.4
125,531.9
126,531.7
\end{lstlisting}

These coordinates are used for 3D reconstruction.

% ==================================================
\section{Camera Calibration}

\textbf{Date: 06 February 2026}

Camera intrinsics are obtained using a checkerboard calibration method.

OpenCV is used to estimate:

\begin{itemize}
  \item Focal lengths $(f_x, f_y)$
  \item Principal point $(c_x, c_y)$
  \item Distortion coefficients
\end{itemize}

Calibration is resolution- and focus-dependent.

Refocusing requires recalibration.

% ==================================================
\section{Checkerboard Calibration Pattern and OpenCV Usage}

\textbf{Date: 09 February 2026}

\subsection{Checkerboard Pattern Requirements}

Camera calibration is performed using a planar black--white checkerboard pattern.

The following parameters are important:

\begin{itemize}
  \item \textbf{Inner corners:} The number of inner corner intersections in the pattern (rows $\times$ columns).
  Typical values used are $9\times6$, $10\times7$, or $11\times8$.
  \item \textbf{Square size:} The physical length of one square edge, measured in millimetres.
  \item \textbf{Overall board size:} The checkerboard should occupy a significant portion of the camera field of view during calibration.
\end{itemize}

The inner corner count and square size must be supplied correctly to the calibration software.

\subsection{Influence of Checkerboard Size}

The absolute square size does not affect intrinsic parameters directly, as these are estimated in pixel units. However, accurate square dimensions are required for:

\begin{itemize}
  \item Estimation of camera pose
  \item Laser plane calibration
  \item Metric reconstruction of 3D geometry
\end{itemize}

For reliable detection, each square should appear approximately 20--80 pixels wide in the captured images.

Boards that are too small produce unreliable corner detection, while boards that are too large limit pose variation.

\subsection{Print Quality and Mounting}

Calibration accuracy is influenced by physical board quality.

Recommended practices include:

\begin{itemize}
  \item High-resolution printing
  \item Matte paper to reduce reflections
  \item Mounting on rigid backing (cardboard or foamboard)
  \item Avoiding warping or bending
\end{itemize}

A flat and rigid checkerboard improves geometric consistency.

\subsection{Use of OpenCV Calibration Functions}

The calibration pipeline is implemented using a custom Python script that interfaces with OpenCV.

The following OpenCV functions perform the core computations:

\begin{itemize}
  \item \texttt{cv2.findChessboardCorners}: Detects checkerboard corner locations
  \item \texttt{cv2.cornerSubPix}: Refines corner positions to subpixel accuracy
  \item \texttt{cv2.calibrateCamera}: Estimates camera intrinsics and distortion
  \item \texttt{cv2.projectPoints}: Used for reprojection error analysis
  \item \texttt{cv2.undistort}: Generates undistorted validation images
\end{itemize}

These functions implement established calibration algorithms and are extensively validated in the computer vision community.

\subsection{Custom Calibration Script Design}

Although OpenCV performs the mathematical optimisation, a custom script has been developed to manage the calibration workflow.

The script provides:

\begin{itemize}
  \item Automated image loading and filtering
  \item Detection success logging
  \item Optional debug image generation
  \item Reprojection error computation
  \item Export of calibration parameters to YAML format
\end{itemize}

This approach combines reliable library functions with a project-specific processing pipeline.

\subsection{Calibration Validity Conditions}

Camera calibration remains valid only under fixed optical conditions.

Recalibration is required if any of the following change:

\begin{itemize}
  \item Lens focus
  \item Image resolution
  \item Digital cropping or scaling
  \item Camera module replacement
\end{itemize}

Changes in camera position or orientation alone do not invalidate intrinsic calibration.

\subsection{Recommended Calibration Practice}

For this project, the following procedure is recommended:

\begin{enumerate}
  \item Mount camera rigidly
  \item Fix focus settings
  \item Select final operating resolution
  \item Capture 15--30 checkerboard images
  \item Vary board distance and orientation
  \item Perform calibration
  \item Verify reprojection error and undistortion
\end{enumerate}

A mean reprojection error below 0.6 pixels is considered acceptable for this application.


% ==================================================
\section{Laser Plane Calibration (Planned)}

\textbf{Date: 07 February 2026}

Laser plane calibration will use a planar checkerboard target.

Proposed method:

\begin{enumerate}
  \item Capture board poses with laser on
  \item Estimate board pose
  \item Back-project laser pixels to rays
  \item Intersect rays with board plane
  \item Fit plane to 3D points
\end{enumerate}

Plane model:

\begin{equation}
ax + by + cz + d = 0
\end{equation}

Parameters will be stored in YAML.

% ==================================================
\section{3D Reconstruction Method}

\textbf{Date: 07 February 2026}

Given:

\begin{itemize}
  \item Camera intrinsics
  \item Laser plane parameters
  \item Pixel centreline
\end{itemize}

Each pixel $(u,v)$ defines a ray:

\begin{equation}
\mathbf{r} = K^{-1}
\begin{bmatrix}
u \\ v \\ 1
\end{bmatrix}
\end{equation}

Intersection with the laser plane yields 3D points.

% ==================================================
\section{Limitations of Current Method}

\textbf{Date: 08 February 2026}

Current limitations:

\begin{itemize}
  \item Sensitive to multiple reflections
  \item Assumes single intersection per row
  \item Limited accuracy for curved stripes
  \item No temporal smoothing
\end{itemize}

The row-mean method is an approximation suitable for early prototyping.

% ==================================================
\section{Future Improvements}

\textbf{Date: 08 February 2026}

Planned upgrades:

\begin{itemize}
  \item Gaussian ridge fitting
  \item Skeletonisation of laser stripe
  \item Robust outlier rejection
  \item Ring laser support
  \item Real-time processing
  \item ML-based defect detection
\end{itemize}

% ==================================================
\section{Development Log}

\subsection*{04 Feb 2026}
Initial codebase created. YAML configuration implemented.

\subsection*{05 Feb 2026}
Laser extraction pipeline operational. HSV thresholding validated.

\subsection*{06 Feb 2026}
PNG input testing successful. CSV output verified.

\subsection*{07 Feb 2026}
Calibration pipeline drafted. Laser plane method planned.

\subsection*{08 Feb 2026}
Limitations of row-wise centroid method identified.

% ==================================================
\section{Conclusion}

This document provides a structured reference for system development and will be expanded as the project progresses.

It supports reproducibility, debugging, and later academic reporting.

\end{document}
